%%
%% This is file `sigconf.tex',
%% generated with the docstrip utility.
%%
%% The original source files were:
%%
%% samples.dtx  (with options: `all,proceedings,sigconf')
%% 
%% IMPORTANT NOTICE:
%% 
%% For the copyright see the source file.
%% 
%% Any modified versions of this file must be renamed
%% with new filenames distinct from sigconf.tex.
%% 
%% For distribution of the original source see the terms
%% for copying and modification in the file samples.dtx.
%% 
%% This generated file may be distributed as long as the
%% original source files, as listed above, are part of the
%% same distribution. (The sources need not necessarily be
%% in the same archive or directory.)
%%
%%
%% Commands for TeXCount
%TC:macro \cite [option:text,text]
%TC:macro \citep [option:text,text]
%TC:macro \citet [option:text,text]
%TC:envir table 0 1
%TC:envir table* 0 1
%TC:envir tabular [ignore] word
%TC:envir displaymath 0 word
%TC:envir math 0 word
%TC:envir comment 0 0
%%
%% The first command in your LaTeX source must be the \documentclass
%% command.
%%
%% For submission and review of your manuscript please change the
%% command to \documentclass[manuscript, screen, review]{acmart}.
%%
%% When submitting camera ready or to TAPS, please change the command
%% to \documentclass[sigconf]{acmart} or whichever template is required
%% for your publication.
%%
%%

\documentclass[sigconf, anonymous, review]{acmart}
\usepackage{algorithm}
\usepackage{algpseudocode}
%%
%% \BibTeX command to typeset BibTeX logo in the docs
\AtBeginDocument{%
  \providecommand\BibTeX{{%
    Bib\TeX}}}

%% Rights management information.  This information is sent to you
%% when you complete the rights form.  These commands have SAMPLE
%% values in them; it is your responsibility as an author to replace
%% the commands and values with those provided to you when you
%% complete the rights form.
\setcopyright{acmlicensed}
\copyrightyear{2026}
\acmYear{2026}
\acmDOI{XXXXXXX.XXXXXXX}
%% These commands are for a PROCEEDINGS abstract or paper.
\acmConference[Conference acronym 'XX]{Make sure to enter the correct
  conference title from your rights confirmation email}{June 03--05,
  2018}{Woodstock, NY}
%%
%%  Uncomment \acmBooktitle if the title of the proceedings is different
%%  from ``Proceedings of ...''!
%%
%%\acmBooktitle{Woodstock '18: ACM Symposium on Neural Gaze Detection,
%%  June 03--05, 2018, Woodstock, NY}
\acmISBN{978-1-4503-XXXX-X/2018/06}


%%
%% Submission ID.
%% Use this when submitting an article to a sponsored event. You'll
%% receive a unique submission ID from the organizers
%% of the event, and this ID should be used as the parameter to this command.
%%\acmSubmissionID{123-A56-BU3}

%%
%% For managing citations, it is recommended to use bibliography
%% files in BibTeX format.
%%
%% You can then either use BibTeX with the ACM-Reference-Format style,
%% or BibLaTeX with the acmnumeric or acmauthoryear sytles, that include
%% support for advanced citation of software artefact from the
%% biblatex-software package, also separately available on CTAN.
%%
%% Look at the sample-*-biblatex.tex files for templates showcasing
%% the biblatex styles.
%%
%%
%% The majority of ACM publications use numbered citations and
%% references.  The command \citestyle{authoryear} switches to the
%% "author year" style.
%%
%% If you are preparing content for an event
%% sponsored by ACM SIGGRAPH, you must use the "author year" style of
%% citations and references.
%% Uncommenting
%% the next command will enable that style.
%%\citestyle{acmauthoryear}


%%
%% end of the preamble, start of the body of the document source.
\begin{document}

%%
%% The "title" command has an optional parameter,
%% allowing the author to define a "short title" to be used in page headers.
\title{FAST-STG: Forecasting-Aligned Stochastic Transport for Probabilistic Urban Spatio-Temporal Graph Forecasting}

%%
%% The "author" command and its associated commands are used to define
%% the authors and their affiliations.
%% Of note is the shared affiliation of the first two authors, and the
%% "authornote" and "authornotemark" commands
%% used to denote shared contribution to the research.
%%
%% The abstract is a short summary of the work to be presented in the
%% article.
\begin{abstract}
Urban sensing underpins Web-facing services such as traffic information
platforms and environmental monitoring, where accurate forecasts and
reliable uncertainty estimates are both important. Existing generative
methods provide flexible tools for probabilistic spatio-temporal graph
forecasting, yet their stochastic formulations may be misaligned with
the conditional structure of forecasting. We identify two such
misalignments: \emph{generation-object misalignment}, where observed history is included in the stochastic state despite already being known, and \emph{generation-coordinate misalignment}, where raw future
values entangle predictable location and scale with stochastic
variation. To address these issues, we propose \textbf{FAST-STG}, a forecasting-aligned stochastic transport framework that keeps the observed history fixed and unperturbed as conditioning and restricts stochastic evolution to the unknown future. FAST-STG further estimates node- and horizon-specific conditional location and scale from the
observed history and spatial graph, and performs transport over
standardized future innovations. The coordinate estimator is pretrained
and frozen before transport learning, providing a stable conditional
coordinate system. Experiments across six real-world datasets spanning air quality and traffic flow demonstrate competitive point forecasting accuracy, while probabilistic evaluations show consistent improvements in probabilistic forecast quality. Code is available at
\url{https://anonymous.4open.science/r/FAST-STG-7D84}.
\end{abstract}

%%
%% The code below is generated by the tool at http://dl.acm.org/ccs.cfm.
%% Please copy and paste the code instead of the example below.
%%
%%
%% Keywords. The author(s) should pick words that accurately describe
%% the work being presented. Separate the keywords with commas.
\keywords{probabilistic forecasting, spatio-temporal graph forecasting,
stochastic transport, uncertainty quantification, urban sensing}
%% A "teaser" image appears between the author and affiliation
%% information and the body of the document, and typically spans the
%% page.
%%
%% This command processes the author and affiliation and title
%% information and builds the first part of the formatted document.
\maketitle

\section{Introduction}
Urban sensing provides an important data layer for Web-facing urban
services, including navigation systems, traffic-information platforms,
environmental dashboards, and decision support
services~\cite{wang2022spatiotemporal,tedjopurnomo2022traffic,
mehmood2017iot}. These infrastructures collect time-series observations
from distributed sensing locations, giving rise to graph-structured
spatio-temporal signals governed by local temporal dynamics and spatial
dependencies~\cite{guo2019attention,bai2020adaptive}. For such Web-connected sensing services, a key challenge
is to transform continuously updated graph-structured observations into
reliable forecasts that characterize both future system states and
their associated uncertainty. Traffic flow
forecasting~\cite{yuan2021traffic,zhang2026msthg} and air quality
forecasting~\cite{zhu2021airquality,liang2023airformer,
han2021jointair,wu2026airdde} are two representative applications.
Despite differences in sensing mechanisms, temporal resolutions, and
spatial dependencies, both require inferring future system states from
historical observations over sensor networks~\cite{wang2020pm25gnn,
zheng2020gman}. As illustrated in Figure~\ref{fig:motivation}, the
observed history and spatial graph provide fixed conditioning
information, whereas the future remains unknown. Conventional point
predictors typically produce a single forecast
trajectory~\cite{yu2018spatio,li2018diffusion,wu2019graph}, providing
limited information about the range and likelihood of plausible future
outcomes~\cite{gneiting2007scoring,salinas2020deepar}. Quantifying this uncertainty is important for Web-connected urban
information services because similar point forecasts may correspond to
substantially different levels of predictive uncertainty. This motivates probabilistic spatio-temporal graph forecasting as a foundation for such services, aiming to characterize the
conditional distribution of future graph signals given historical
observations and spatial dependencies~\cite{rasul2021timegrad,
wen2023diffstg}.

\begin{figure}[t]
    \centering
    \includegraphics[width=\columnwidth]
    {Figs/Fig1_Motivation.pdf}
    \caption{
        Motivation and application setting of probabilistic forecasting
        for Web-connected urban sensing. Historical time-series observations
        and spatial graph structures provide deterministic context, whereas
        unknown future graph signals are modeled probabilistically. The
        illustrated platform is conceptual rather than an implemented
        online system.
    }
    \label{fig:motivation}
\end{figure}

Probabilistic forecasting has progressed from likelihood-based sequence
models~\cite{salinas2020deepar} toward more flexible generative
approaches~\cite{rasul2021timegrad,tashiro2021csdi,kollovieh2025tsflow}. Autoregressive
methods such as DeepAR~\cite{salinas2020deepar} parameterize predictive
distributions with tractable likelihoods and recursively generate future
observations. Diffusion-based methods such as
TimeGrad~\cite{rasul2021timegrad} and
CSDI~\cite{tashiro2021csdi} introduce iterative denoising into
probabilistic time-series forecasting and imputation, while
DiffSTG~\cite{wen2023diffstg} extends this paradigm to spatio-temporal
graphs by incorporating graph-temporal dependencies into the denoising
process. More recently, flow matching~\cite{lipman2023flow} and
rectified flow~\cite{liu2023rectified} formulate generation as
continuous transport between distributions, with
TSFlow~\cite{kollovieh2025tsflow} adapts conditional flow matching to
probabilistic time-series forecasting. Together, these developments
have expanded probabilistic forecasting from explicit distribution
parameterization to diffusion and continuous transport. However,
existing studies have largely focused on \emph{how} flexible predictive
distributions are generated, while the forecasting-specific question
of \emph{which variables} should participate in stochastic generation
has received relatively limited attention.

This unresolved issue reveals a fundamental asymmetry in conditional
forecasting: the historical window $\mathbf{X}_H$ is observed
conditioning, whereas the future window $\mathbf{X}_F$ remains the
unknown prediction target, as reflected in likelihood-based, diffusion,
and flow-based forecasting formulations~\cite{salinas2020deepar,
rasul2021timegrad,kollovieh2025tsflow}. Conditional generative models~\cite{tashiro2021csdi,liu2023pristi}
typically distinguish observed variables from those to be inferred
through masking or conditioning mechanisms, while a recent study~\cite{li2026backtobasics} further shows that the choice of prediction target can materially affect generative
learning. Diffusion~\cite{ho2020ddpm} and
score-based~\cite{song2021score} formulations apply stochastic
perturbation and recovery to variables included in the generative state.
If this state extends beyond the unknown future, observed history may
also undergo unnecessary stochastic evolution and reconstruction.
This occurs in generalized STG diffusion formulations such as
DiffSTG~\cite{wen2023diffstg}, where history and future are placed
within a shared diffusion state. We term this structural discrepancy
\emph{generation-object misalignment}. Accordingly, a
forecasting-aligned generative process should preserve
$\mathbf{X}_H$ as fixed conditioning and restrict stochastic evolution
to $\mathbf{X}_F$.

Once stochastic evolution is restricted to the unknown future, a second
question arises: in which coordinates should that future be modeled? A recent generative study~\cite{li2026backtobasics} suggests that target
parameterization can substantially affect the learning problem,
while forecasting methods have increasingly separated predictable
structure from stochastic variation. Re-Diffusion~\cite{zhang2026rediffusion} models residual errors from a pretrained forecaster, and
NsDiff~\cite{ye2025nsdiff} incorporates pretrained conditional mean and variance estimates
through a location and scale noise model. CW-Gen~\cite{yang2026cwgen} whitens forecasting targets using estimated conditional means and
sliding-window covariances for both diffusion and flow matching, while CoST~\cite{sheng2026cost} combines conditional mean prediction
with residual diffusion and spatially adaptive scale modeling. These developments indicate the benefit of
separating predictable conditional structure from stochastic variation.
However, in spatio-temporal graph forecasting, conditional location and
uncertainty scale can vary across both nodes and forecast horizons
according to the observed history and graph. Directly transporting the
raw future therefore requires the generative field to jointly represent
predictable location, heterogeneous scale, and residual distributional
shape. We refer to this discrepancy as
\emph{generation-coordinate misalignment}, where the stochastic target
remains expressed in coordinates entangled with context-dependent
location and scale.

To address both forms of misalignment, we propose
\textbf{FAST-STG}, a \textbf{f}orecasting-\textbf{a}ligned
\textbf{s}tochastic \textbf{t}ransport framework for probabilistic
spatio-temporal graph forecasting. FAST-STG keeps the observed history
fixed and unperturbed as conditioning and restricts stochastic evolution
to the unknown future, thereby aligning the stochastic state with the
forecasting task. It further estimates node- and horizon-specific
conditional location and scale from the observed history and spatial
graph, and performs transport over standardized future innovations
rather than raw future values. The coordinate estimator is pretrained
and frozen before transport learning, providing a stable conditional
coordinate system throughout transport optimization. Through this
design, FAST-STG keeps observed contextual evidence outside the
stochastic state and separates predictable conditional structure from
stochastic variation.

Our main contributions are summarized as follows:

\begin{itemize}

\item We identify two forms of misalignment between generative modeling
and spatio-temporal forecasting:
\emph{generation-object misalignment}, where stochastic evolution
extends beyond the unknown future variables, and
\emph{generation-coordinate misalignment}, where the stochastic target
remains expressed in coordinates entangled with conditional location
and scale. This formulation highlights two forecasting-specific design
principles for stochastic generation.

\item We propose FAST-STG, which restricts stochastic transport to the
unknown future and constructs node- and horizon-specific conditional
location and scale coordinates from the observed history and spatial
graph. FAST-STG transports standardized future innovations and employs
a staged training strategy that freezes the coordinate estimator before
transport learning, thereby separating predictable conditional
structure from stochastic variation.

\item We evaluate FAST-STG on six real-world datasets from two urban
sensing domains, including four traffic flow datasets and two air
quality datasets. Experiments cover point forecasting accuracy,
probabilistic forecast quality, component ablations, sensitivity
analysis, qualitative forecasting behavior, and computational
efficiency, demonstrating the effectiveness of FAST-STG across both
domains.

\end{itemize}

\section{Related Work}
\subsection{Probabilistic and Generative Spatio-temporal Forecasting}

Existing spatio-temporal forecasting methods have developed increasingly
expressive mechanisms for modeling dependencies over sensor networks.
STGCN~\cite{yu2018spatio} combines graph convolution with temporal
convolution, while DCRNN~\cite{li2018diffusion} integrates diffusion
convolution into recurrent sequence modeling. Graph
WaveNet~\cite{wu2019graph} and AGCRN~\cite{bai2020adaptive} further
learn adaptive spatial dependencies from data. More recent approaches
extend graph forecasting through attention and Transformer
architectures~\cite{liang2023airformer,han2021jointair}, continuous
dynamics~\cite{long2024delay,wu2026airdde}, and hypergraph
modeling~\cite{zhang2026msthg}. These methods substantially improve
spatio-temporal representation learning, but primarily focus on point
forecasting and provide limited characterization of predictive
uncertainty.

Probabilistic forecasting addresses this limitation through explicit
distribution modeling or stochastic generation. DeepAR
~\cite{salinas2020deepar} parameterizes predictive distributions with
autoregressive likelihoods. TimeGrad~\cite{rasul2021timegrad} introduces
conditional diffusion for probabilistic forecasting, while
CSDI~\cite{tashiro2021csdi} learns conditional distributions of
unobserved values. Flow matching~\cite{lipman2023flow}, rectified
flow~\cite{liu2023rectified}, and stochastic
interpolants~\cite{albergo2025stochastic} provide alternative
distribution transport formulations, with TSFlow
~\cite{kollovieh2025tsflow} adapting conditional flow matching to
probabilistic time-series forecasting. For graph-structured data,
DiffSTG~\cite{wen2023diffstg} incorporates spatio-temporal dependencies
into denoising diffusion, GCRDD~\cite{li2023gcrdd} combines graph
recurrent modeling with conditional diffusion, and
USTD~\cite{hu2024unifying} develops a unified diffusion formulation for
probabilistic spatio-temporal graph learning. SpecSTG
~\cite{lin2025specstg} further performs probabilistic generation in the
spectral domain. Despite their different generative mechanisms, these
methods mainly advance distributional flexibility and spatio-temporal
dependency modeling, while the explicit alignment between the
stochastic state and the conditional structure of forecasting remains
comparatively underexplored.

\subsection{Residual and Conditional Coordinate Modeling}

A growing line of probabilistic forecasting research reduces the burden
of generative modeling by separating predictable structure from
stochastic variation. Re-Diffusion~\cite{zhang2026rediffusion} applies a
pretrained forecaster and models the resulting residual distribution in
latent space, thereby concentrating generative learning on deviations
that remain unexplained by the deterministic prediction.
CoST~\cite{sheng2026cost} similarly combines conditional mean prediction
with residual diffusion and further accounts for spatially varying
uncertainty scales. Related hybrid approaches decompose the forecasting
target according to uncertainty or temporal structure.
D$^3$U~\cite{li2025d3u} assigns high-certainty components to a point
forecaster and models uncertain components with conditional diffusion,
while TCDM~\cite{xu2025tcdm} separates trend and residual dynamics
before applying diffusion to the residual component. HySRD-Net
~\cite{chen2026hysrd} extends trend-residual decomposition to
probabilistic spatio-temporal graph forecasting and models residual
uncertainty in the graph spectral domain. Collectively, these methods
show that removing predictable components can simplify the stochastic
learning problem and allow generative models to focus on uncertainty
that remains after deterministic structure has been extracted. Their
residual representations, however, are largely determined by the
preceding decomposition or point predictor, and do not necessarily
provide an explicit conditional coordinate system that adapts jointly
to spatial location and forecast horizon.

A complementary direction explicitly adapts the coordinates in which
the stochastic target is modeled. NsDiff~\cite{ye2025nsdiff} constructs
a location and scale noise model from pretrained conditional mean and
variance estimates, allowing diffusion to operate relative to the
conditional statistics of the future rather than directly in the raw
observation space. CW-Gen~\cite{yang2026cwgen} further whitens
forecasting targets using estimated conditional means and covariances
and applies the resulting representation to both diffusion and flow
matching. Compared with residual decomposition alone, such approaches
more directly modify the geometry and scale of the generative target,
highlighting the importance of conditional reparameterization for
distribution learning. Nevertheless, spatio-temporal graph forecasting
introduces additional heterogeneity because conditional location and
uncertainty scale can vary simultaneously across sensing nodes and
forecast horizons according to both historical states and spatial
dependencies. The joint construction of graph-conditioned,
node- and horizon-specific coordinates, together with a staged training
strategy that keeps these coordinates fixed during subsequent transport
learning, remains comparatively underexplored.

\section{Methodology}
\label{sec:method}

This section presents FAST-STG for probabilistic spatio-temporal graph
forecasting. FAST-STG aligns stochastic generation with forecasting
along two dimensions: it restricts stochastic evolution to the unknown
future and constructs conditional location and scale coordinates for
future transport. We first formulate the forecasting problem and then
describe conditional coordinate pretraining, future-only transport
learning, and stochastic sampling for probabilistic prediction.

\subsection{Problem Formulation}
\label{sec:problem_formulation}

\begin{figure*}[t]
    \centering
    \includegraphics[width=0.90\textwidth]{Figs/Fig2.pdf}
    \caption{
    Overall framework of FAST-STG.
    (a) Conditional coordinate pretraining estimates node- and
    horizon-specific conditional location and scale from historical
    observations and graph structure.
    (b) Future-only transport training learns stochastic transport in
    the standardized future innovation space while keeping historical
    observations fixed as conditioning.
    }
    \label{fig:framework}
\end{figure*}

Let $\mathcal{G}=(\mathcal{V},\mathcal{E},\mathbf{A})$ denote a
spatial graph, where $\mathcal{V}$ contains $N$ sensor nodes,
$\mathcal{E}$ represents their spatial relations, and
$\mathbf{A}\in\mathbb{R}^{N\times N}$ is the adjacency matrix. At time
step $t$, observations over all nodes are denoted by
$\mathbf{x}_{t}\in\mathbb{R}^{N\times F}$, where $F$ is the number of
observed features per node. Given a historical window of length
$T_h$ and a forecast horizon of length $T_p$, we define
%
\begin{equation}
\begin{aligned}
\mathbf{X}^{\mathrm{his}}
&=
[\mathbf{x}_{t-T_h+1},\ldots,\mathbf{x}_{t}]
\in\mathbb{R}^{T_h\times N\times F},\\
\mathbf{X}^{\mathrm{fut}}
&=
[\mathbf{x}_{t+1},\ldots,\mathbf{x}_{t+T_p}]
\in\mathbb{R}^{T_p\times N\times F}.
\end{aligned}
\label{eq:history_future}
\end{equation}
%
Probabilistic spatio-temporal graph forecasting aims to characterize
the conditional distribution of the unknown future given the observed
history and spatial graph:
%
\begin{equation}
p\left(
\mathbf{X}^{\mathrm{fut}}
\mid
\mathbf{X}^{\mathrm{his}},\mathbf{A}
\right).
\label{eq:forecasting_distribution}
\end{equation}
%
At inference time, this distribution is represented by $S$ sampled
future trajectories,
$\{\widehat{\mathbf{X}}^{\mathrm{fut},(s)}\}_{s=1}^{S}$,
which characterize plausible future realizations and predictive
uncertainty.

\subsection{Conditional Coordinate Pretraining}
\label{sec:coordinate_pretraining}

FAST-STG first learns a conditional coordinate system before stochastic
transport. As illustrated in Figure~\ref{fig:framework}(a), a
graph-conditioned location and scale estimator $f_{\phi}$ takes the
observed history and spatial graph as input and predicts node- and
horizon-specific coordinate parameters:
%
\begin{equation}
\left(
\boldsymbol{\mu}_{\phi},
\boldsymbol{\sigma}_{\phi}
\right)
=
f_{\phi}
\left(
\mathbf{X}^{\mathrm{his}},
\mathbf{A}
\right),
\label{eq:coordinate_estimator}
\end{equation}
%
where
$\boldsymbol{\mu}_{\phi},
\boldsymbol{\sigma}_{\phi}
\in\mathbb{R}^{T_p\times N\times F}$.

To capture temporal context and spatial dependencies, the estimator
first encodes the historical observations and then propagates the
resulting representations over the sensor graph:
%
\begin{equation}
\mathbf{H}
=
\operatorname{Enc}_{\phi}
\left(
\mathbf{X}^{\mathrm{his}}
\right),
\qquad
\mathbf{G}
=
\operatorname{GraphProp}_{\phi}
\left(
\mathbf{H},\mathbf{A}
\right),
\label{eq:coordinate_encoding}
\end{equation}
%
where $\mathbf{G}=\{\mathbf{g}_{v}\}_{v=1}^{N}$ contains a
graph-aware representation for each sensor node. For each forecast
horizon $j\in\{1,\ldots,T_p\}$, corresponding to the future time step
$t+j$, we introduce a learnable horizon embedding $\mathbf{e}_{j}$ and
construct
\begin{equation}
\mathbf{q}_{v,j}
=
\operatorname{Dec}_{\phi}
\left(
\mathbf{g}_{v}
+
\mathbf{e}_{j}
\right).
\label{eq:coordinate_horizon_representation}
\end{equation}
The conditional location is parameterized as a correction to the most
recent observation, while the conditional scale is bounded to remain
positive:
\begin{equation}
\begin{gathered}
\boldsymbol{\mu}_{\phi,v,j}
=
\mathbf{x}_{t,v}
+
\alpha_{\mu}
\left(
\mathbf{W}_{\mu}\mathbf{q}_{v,j}
+
\mathbf{b}_{\mu}
\right),
\\
\boldsymbol{\sigma}_{\phi,v,j}
=
\sigma_{\min}
+
\left(
\sigma_{\max}-\sigma_{\min}
\right)
\operatorname{sigmoid}
\left(
\mathbf{W}_{\sigma}\mathbf{q}_{v,j}
+
\mathbf{b}_{\sigma}
\right).
\end{gathered}
\label{eq:conditional_location_scale}
\end{equation}
Here, $\alpha_{\mu}$ controls the magnitude of the location correction,
while $\sigma_{\min}$ and $\sigma_{\max}$ specify the lower and upper
bounds of the conditional scale. This parameterization allows the
coordinate system to adapt across sensor nodes and forecast horizons
according to the observed history and spatial graph.

The coordinate estimator is pretrained using the future observations
$\mathbf{X}^{\mathrm{fut}}$ by minimizing the following diagonal
Gaussian negative log likelihood, denoted as the location and scale
loss $\mathcal{L}_{\mathrm{LS}}$:
%
\begin{equation}
\mathcal{L}_{\mathrm{LS}}(\phi)
=
\operatorname{mean}
\left[
\frac{1}{2}
\left(
\frac{
\mathbf{X}^{\mathrm{fut}}
-
\boldsymbol{\mu}_{\phi}
}{
\boldsymbol{\sigma}_{\phi}
}
\right)^{2}
+
\log
\boldsymbol{\sigma}_{\phi}
\right],
\label{eq:ls_loss}
\end{equation}
%
where all operations are applied element-wise. This objective is used
only to estimate the conditional location and scale coordinates and
does not impose a Gaussian form on the final predictive distribution.
After pretraining, the estimator parameters $\phi$ are frozen, yielding
a fixed conditional coordinate system for subsequent future-only
transport learning.

\subsection{Future-Only Transport Training}
\label{sec:future_transport_training}

After conditional coordinate pretraining, the estimator parameters
$\phi$ are frozen and the learned location and scale define the
coordinate system for transport. The future observations are
standardized as
%
\begin{equation}
\mathbf{R}^{\mathrm{fut}}
=
\left(
\mathbf{X}^{\mathrm{fut}}
-
\boldsymbol{\mu}_{\phi}
\right)
\oslash
\boldsymbol{\sigma}_{\phi},
\label{eq:standardized_innovation}
\end{equation}
%
where $\oslash$ denotes element-wise division. The resulting
$\mathbf{R}^{\mathrm{fut}}
\in\mathbb{R}^{T_p\times N\times F}$
represents the standardized future innovation obtained by centering
on the conditional location and normalizing by the conditional scale. Because the coordinate estimator remains frozen, this target
representation is fixed throughout transport learning.

FAST-STG performs stochastic transport only over the unknown future
while keeping the observed history fixed and unperturbed as conditioning. We first
sample a Gaussian source
%
\begin{equation}
\mathbf{Z}_{0}
\sim
\mathcal{N}
\left(
\mathbf{0},
\mathbf{I}
\right),
\label{eq:transport_source}
\end{equation}
%
where
$\mathbf{Z}_{0}\in\mathbb{R}^{T_p\times N\times F}$,
and draw a continuous transport time
$\tau\sim\mathcal{U}(0,1)$. We then construct a linear path between the
Gaussian source and the standardized future innovation:
%
\begin{equation}
\mathbf{Z}_{\tau}
=
(1-\tau)\mathbf{Z}_{0}
+
\tau\mathbf{R}^{\mathrm{fut}}.
\label{eq:linear_transport_path}
\end{equation}
%
The stochastic state therefore remains defined exclusively over the
future interval throughout transport learning.

To provide the graph-temporal backbone with complete temporal context,
we concatenate the evolving future state with the fixed historical
observations and construct a masked conditioning sequence:
%
\begin{equation}
\overline{\mathbf{Z}}_{\tau}
=
\operatorname{Cat}_{t}
\left(
\mathbf{X}^{\mathrm{his}},
\mathbf{Z}_{\tau}
\right),
\qquad
\mathbf{C}
=
\operatorname{Cat}_{t}
\left(
\mathbf{X}^{\mathrm{his}},
\mathbf{0}
\right),
\label{eq:future_only_inputs}
\end{equation}
%
where $\operatorname{Cat}_{t}$ denotes concatenation along the temporal
dimension and
$\mathbf{0}\in\mathbb{R}^{T_p\times N\times F}$
masks the unobserved future interval. Although
$\overline{\mathbf{Z}}_{\tau}$ spans the complete historical and future
window, its historical segment is copied directly from
$\mathbf{X}^{\mathrm{his}}$ and remains unchanged. Hence, only
$\mathbf{Z}_{\tau}$ participates in stochastic evolution.

Following DiffSTG~\cite{wen2023diffstg}, we employ UGNet as the
graph-temporal backbone, allowing the comparison to focus on the
stochastic formulation rather than changes to the backbone
architecture. Given the state-augmented sequence, masked condition,
transport time, and spatial graph, UGNet predicts the clean future
innovation:
%
\begin{equation}
\widehat{\mathbf{R}}_{\theta}^{\mathrm{fut}}
=
\Pi_{\mathrm{fut}}
\left[
\operatorname{UGNet}_{\theta}
\left(
\overline{\mathbf{Z}}_{\tau},
\tau;
\mathbf{C},
\mathbf{A}
\right)
\right],
\label{eq:future_only_prediction}
\end{equation}
%
where $\Pi_{\mathrm{fut}}$ retains the final $T_p$ time steps of the
backbone output. The historical output is discarded and does not enter
the transport objective. As illustrated in
Figure~\ref{fig:ugnet}, UGNet consists of stacked spatio-temporal
residual blocks that capture temporal and spatial dependencies through
temporal convolution and graph convolution operations.

\begin{figure}[t]
    \centering
    \includegraphics[width=\columnwidth]{Figs/UGNet-Final.pdf}
    \caption{
    Architecture of the conditioned UGNet backbone used for future
    innovation prediction.
    }
    \label{fig:ugnet}
\end{figure}

Rather than directly supervising a transport velocity, FAST-STG learns
the clean endpoint using the following transport loss
$\mathcal{L}_{\mathrm{T}}$:
%
\begin{equation}
\mathcal{L}_{\mathrm{T}}(\theta)
=
\mathbb{E}
\left[
\left\|
\widehat{\mathbf{R}}_{\theta}^{\mathrm{fut}}
-
\mathbf{R}^{\mathrm{fut}}
\right\|_{2}^{2}
\right],
\label{eq:transport_loss}
\end{equation}
%
where the expectation is taken over training samples, Gaussian source
noise, and transport times. During this stage, gradients update only
the UGNet parameters $\theta$, while the pretrained coordinate
estimator remains frozen. The resulting model learns future-only
transport in a fixed conditional innovation space.

\subsection{Stochastic Sampling and Inverse Transformation}
\label{sec:sampling_reconstruction}

\begin{figure}[t]
    \centering
    \includegraphics[width=\columnwidth]{Figs/SDE采样示意图.pdf}
    \caption{
    Stochastic sampling procedure of FAST-STG. Starting from a Gaussian
    future innovation state, an endpoint-induced conditional transport
    field is integrated through $K$ Euler--Maruyama steps while the
    observed history remains fixed as conditioning. Repeating the
    procedure with independent source and noise samples yields multiple
    probabilistic forecasts.
    }
    \label{fig:sampling}
\end{figure}

At inference time, the pretrained coordinate estimator and transport
backbone are fixed. FAST-STG generates future innovations by integrating
an endpoint-induced conditional transport field, as illustrated in
Figure~\ref{fig:sampling}. For the linear path in
Eq.~\eqref{eq:linear_transport_path},
%
\begin{equation}
\mathbf{R}^{\mathrm{fut}}
-
\mathbf{Z}_{\tau}
=
(1-\tau)
\left(
\mathbf{R}^{\mathrm{fut}}
-
\mathbf{Z}_{0}
\right).
\label{eq:endpoint_velocity_relation}
\end{equation}
%
This relation allows the transport direction to be constructed from the
predicted clean endpoint. We discretize the transport interval into
$K$ sampling steps with $\Delta\tau=1/K$ and
$\tau_k=k\Delta\tau$ for $k=0,\ldots,K-1$. At sampling step $k$,
UGNet re-estimates the clean future innovation from the current state
$\mathbf{Z}_{k}$ under the same fixed historical conditioning used
during training. Denoting this prediction by
$\widehat{\mathbf{R}}_{\theta,k}^{\mathrm{fut}}$, we define
%
\begin{equation}
\mathbf{v}_{\theta}
\left(
\mathbf{Z}_{k},\tau_k
\mid
\mathbf{X}^{\mathrm{his}},\mathbf{A}
\right)
=
\frac{
\widehat{\mathbf{R}}_{\theta,k}^{\mathrm{fut}}
-
\mathbf{Z}_{k}
}{
1-\tau_k+\delta
},
\label{eq:transport_velocity}
\end{equation}
%
where $\delta>0$ is a small constant for numerical stability.

To maintain stochasticity during generation, FAST-STG augments the
endpoint-induced transport direction with a time-dependent diffusion
term:
%
\begin{equation}
\mathrm{d}\mathbf{Z}_{\tau}
=
\mathbf{v}_{\theta}
\left(
\mathbf{Z}_{\tau},\tau
\mid
\mathbf{X}^{\mathrm{his}},\mathbf{A}
\right)
\mathrm{d}\tau
+
g(\tau)\mathrm{d}\mathbf{W}_{\tau},
\qquad
g(\tau)=\gamma(1-\tau),
\label{eq:transport_sde}
\end{equation}
%
where $\mathbf{W}_{\tau}$ is a standard Wiener process and $\gamma$
controls the stochastic intensity. The diffusion magnitude decreases
as the trajectory approaches the terminal future innovation.

Starting from the Gaussian source in
Eq.~\eqref{eq:transport_source}, the sampling procedure proceeds over
the $K$ discretized transport steps. At each step, the current future
state is concatenated with the unchanged history,
%
\begin{equation}
\overline{\mathbf{Z}}_{k}
=
\operatorname{Cat}_{t}
\left(
\mathbf{X}^{\mathrm{his}},
\mathbf{Z}_{k}
\right),
\label{eq:sampling_input}
\end{equation}
%
and UGNet predicts
$\widehat{\mathbf{R}}_{\theta,k}^{\mathrm{fut}}$
using the fixed masked condition $\mathbf{C}$ defined in
Eq.~\eqref{eq:future_only_inputs}. The future state is then updated by
Euler--Maruyama:
%
\begin{equation}
\mathbf{Z}_{k+1}
=
\mathbf{Z}_{k}
+
\Delta\tau\,
\mathbf{v}_{\theta}
\left(
\mathbf{Z}_{k},\tau_k
\mid
\mathbf{X}^{\mathrm{his}},\mathbf{A}
\right)
+
\gamma(1-\tau_k)
\sqrt{\Delta\tau}\,
\boldsymbol{\epsilon}_{k},
\label{eq:euler_maruyama}
\end{equation}
%
where
$\boldsymbol{\epsilon}_{k}
\sim
\mathcal{N}(\mathbf{0},\mathbf{I})$
is independently sampled at each step. Only the future state evolves
during sampling, while the historical observations remain unchanged.

After the final step, $\mathbf{Z}_{K}$ represents a generated
standardized future innovation. We map it back to the observation space
using the frozen conditional coordinates:
%
\begin{equation}
\widehat{\mathbf{X}}^{\mathrm{fut}}
=
\boldsymbol{\mu}_{\phi}
+
\boldsymbol{\sigma}_{\phi}
\odot
\mathbf{Z}_{K},
\label{eq:inverse_coordinate_transform}
\end{equation}
%
where $\odot$ denotes element-wise multiplication. Repeating the
sampling procedure with $S$ independent Gaussian source and Wiener
noise realizations produces
%
\begin{equation}
\left\{
\widehat{\mathbf{X}}^{\mathrm{fut},(s)}
\right\}_{s=1}^{S},
\label{eq:forecast_samples}
\end{equation}
%
which approximate the conditional predictive distribution.

\section{Experiments}
\label{sec:experiments}

We evaluate FAST-STG on air quality and traffic flow forecasting tasks
from three perspectives. Specifically, we address the following
questions: \textbf{(Q1)} How does FAST-STG compare with existing
probabilistic forecasting methods in terms of predictive distribution
quality and point forecasting accuracy? \textbf{(Q2)} How do
future-only stochastic transport and conditional location and scale
coordinates individually and jointly affect forecasting performance?
\textbf{(Q3)} How sensitive is FAST-STG to its transport and sampling
configurations, and what computational cost does stochastic inference
introduce? The main text reports probabilistic comparisons and
controlled ablations on AIR-BJ, AIR-GZ, and PEMS08, sensitivity
analysis on AIR-GZ, and computational efficiency on AIR-BJ. Additional
point forecasting results on the traffic flow datasets are provided in
Appendix~\ref{app:point_results}.

\subsection{Experimental Setup}
\label{sec:experimental_setup}

\subsubsection{Datasets}

We evaluate FAST-STG on six real-world spatio-temporal datasets from
two urban sensing domains. AIR-BJ and AIR-GZ are used for air quality
forecasting~\cite{wen2023diffstg}, while PEMS03, PEMS04, PEMS07, and
PEMS08 are used for traffic flow forecasting~\cite{zhang2026msthg}.
For all datasets, the first feature is selected as the prediction
target. Dataset statistics are summarized in
Table~\ref{tab:dataset_statistics}.

\begin{table}[h]
    \centering
    \caption{Statistics of the six urban sensing datasets.}
    \label{tab:dataset_statistics}
    \small
    \setlength{\tabcolsep}{5pt}
    \begin{tabular}{@{}llrrc@{}}
        \toprule
        Domain & Dataset & Time steps & Nodes & Interval \\
        \midrule
        Air quality  & AIR-BJ  & 8,760  & 34  & 1 hour \\
        Air quality  & AIR-GZ  & 8,760  & 41  & 1 hour \\
        Traffic flow & PEMS03  & 26,208 & 358 & 5 min  \\
        Traffic flow & PEMS04  & 16,992 & 307 & 5 min  \\
        Traffic flow & PEMS07  & 28,224 & 883 & 5 min  \\
        Traffic flow & PEMS08  & 17,856 & 170 & 5 min  \\
        \bottomrule
    \end{tabular}
\end{table}

\begin{table*}[t]
\centering
\caption{Forecasting performance on AIR-BJ, AIR-GZ, and PEMS08.
Best results are shown in
bold and second-best results are underlined. Methods marked with
$\dagger$ are reproduced under the same experimental protocol as
FAST-STG; the remaining baseline results follow the published
DiffSTG evaluation setting.}
\label{tab:probabilistic_results}
\small
\setlength{\tabcolsep}{3.8pt}
\renewcommand{\arraystretch}{1.08}
\begin{tabular}{@{}llccccccccc@{}}
\toprule
& &
\multicolumn{3}{c}{AIR-BJ}
& \multicolumn{3}{c}{AIR-GZ}
& \multicolumn{3}{c}{PEMS08} \\
\cmidrule(lr){3-5}
\cmidrule(lr){6-8}
\cmidrule(lr){9-11}

Method & Venue
& MAE $\downarrow$ & RMSE $\downarrow$ & CRPS $\downarrow$
& MAE $\downarrow$ & RMSE $\downarrow$ & CRPS $\downarrow$
& MAE $\downarrow$ & RMSE $\downarrow$ & CRPS $\downarrow$ \\
\midrule

Latent ODE~\cite{rubanova2019latent}
& NeurIPS'19
& 20.61 & 32.27 & 0.47
& 12.92 & 18.76 & 0.30
& 26.05 & 39.50 & 0.11 \\

DeepAR~\cite{salinas2020deepar}
& IJF'20
& 20.15 & 32.09 & 0.37
& 11.77 & 17.45 & 0.23
& 21.56 & 33.37 & 0.07 \\

CSDI~\cite{tashiro2021csdi}
& NeurIPS'21
& 26.52 & 40.33 & 0.50
& 13.75 & 19.40 & 0.28
& 32.11 & 47.40 & 0.11 \\

TimeGrad~\cite{rasul2021timegrad}
& ICML'21
& 18.64 & 31.86 & 0.36
& 12.36 & 18.15 & 0.25
& 24.46 & 38.06 & 0.09 \\

MC Dropout~\cite{wu2021uncertainty}
& KDD'21
& 20.80 & 40.54 & 0.45
& 11.12 & 17.07 & 0.25
& 19.01 & 29.35 & 0.07 \\

DiffSTG~\cite{wen2023diffstg}
& SIGSPATIAL'23
& 17.88 & 29.60 & 0.34
& 10.95 & 16.66 & 0.22
& 17.68 & 27.13 & \underline{0.06} \\

SpecSTG$\dagger$~\cite{lin2025specstg}
& ICLR-W'25
& 19.09 & 31.78 & 0.35
& 10.89 & \underline{16.12} & 0.20
& 17.21 & 26.49 & \underline{0.06} \\

NsDiff$\dagger$~\cite{ye2025nsdiff}
& ICML'25
& 19.06 & 30.96 & 0.36
& 11.22 & 16.48 & 0.21
& 20.09 & 30.25 & 0.07 \\

CW-Diff$\dagger$~\cite{yang2026cwgen}
& ICLR'26
& \underline{17.73} & \textbf{25.09} & 0.36
& 24.60 & 38.06 & 0.49
& 25.84 & 39.20 & 0.09 \\

CW-Flow$\dagger$~\cite{yang2026cwgen}
& ICLR'26
& 18.12 & 30.78 & 0.34
& 11.57 & 17.07 & 0.21
& 18.11 & 27.38 & 0.07 \\

CoST$\dagger$~\cite{sheng2026cost}
& KDD'26
& 18.60 & 30.28 & \underline{0.33}
& \underline{10.67} & 16.49 & \underline{0.19}
& \underline{15.89} & \underline{24.98} & \underline{0.06} \\

\midrule

\textbf{FAST-STG}
& \textbf{Ours}
& \textbf{16.35} & \underline{27.57} & \textbf{0.29}
& \textbf{10.00} & \textbf{15.00} & \textbf{0.18}
& \textbf{15.66} & \textbf{24.53} & \textbf{0.05} \\

\bottomrule
\end{tabular}
\end{table*}

\subsubsection{Baselines}
\label{sec:baselines}

For probabilistic forecasting, we compare FAST-STG with eleven
representative baselines. DeepAR~\cite{salinas2020deepar} models
predictive distributions through autoregressive likelihoods, while
Latent ODE~\cite{rubanova2019latent} and MC
Dropout~\cite{wu2021uncertainty} provide continuous-dynamics and
uncertainty-estimation baselines. TimeGrad~\cite{rasul2021timegrad} and
CSDI~\cite{tashiro2021csdi} represent diffusion-based probabilistic
time-series models. For spatio-temporal graph forecasting, we include
DiffSTG~\cite{wen2023diffstg} and SpecSTG~\cite{lin2025specstg}.
Recent approaches that separate predictable conditional structure from
stochastic variation include NsDiff~\cite{ye2025nsdiff},
CW-Gen~\cite{yang2026cwgen}, instantiated as CW-Diff and CW-Flow, and
CoST~\cite{sheng2026cost}. SpecSTG, NsDiff, CW-Diff, CW-Flow, and CoST
are reproduced under the same experimental setting as FAST-STG,
including identical data splits and evaluation protocols.

For point forecasting, we use the mean of the generated trajectories
as the prediction of FAST-STG and compare it with representative
deterministic spatio-temporal forecasting methods, including
STGCN~\cite{yu2018spatio}, DCRNN~\cite{li2018diffusion},
Graph WaveNet~\cite{wu2019graph}, ASTGCN~\cite{guo2019attention}, AGCRN~\cite{bai2020adaptive}, STGODE~\cite{fang2021stgode}, and
STG-NCDE~\cite{choi2022stgncde}. These baselines cover graph
convolution and recurrent modeling, adaptive graph learning,
attention-based forecasting, and continuous spatio-temporal dynamics.
Complete point forecasting results on PEMS03, PEMS04, PEMS07, and
PEMS08 are reported in Appendix~\ref{app:point_results}.

\subsubsection{Evaluation Metrics}
\label{sec:metrics}

We evaluate both point forecasting accuracy and probabilistic forecast
quality. Given $S$ generated trajectories
$\{\hat{y}_{i}^{(s)}\}_{s=1}^{S}$ for target $y_i$, we use their
predictive mean
%
\begin{equation}
\bar{y}_{i}
=
\frac{1}{S}
\sum_{s=1}^{S}
\hat{y}_{i}^{(s)}
\end{equation}
%
as the point forecast and report Mean Absolute Error (MAE) and Root
Mean Squared Error (RMSE), where lower values indicate better point
forecasting accuracy.

For probabilistic evaluation, we use the Continuous Ranked Probability
Score (CRPS)~\cite{gneiting2007scoring} as the primary metric. Let
$F_i$ denote the predictive cumulative distribution for $y_i$. CRPS is
defined as
%
\begin{equation}
\mathrm{CRPS}(F_i,y_i)
=
\int_{\mathbb{R}}
\left(
F_i(z)-\mathbb{I}\{y_i\leq z\}
\right)^{2}
\mathrm{d}z,
\label{eq:crps}
\end{equation}
%
where lower values indicate better probabilistic forecasts. Since
FAST-STG represents the predictive distribution through generated
trajectories, CRPS is estimated empirically from the forecast samples
following the adopted evaluation protocol. We further evaluate
prediction-interval calibration and sharpness using
PICP$_{90}$, PIW$_{90}$, and MIS$_{90}$, with their definitions and
results reported in Appendix~\ref{app:interval_metrics}.

\subsubsection{Implementation Details}
\label{sec:implementation}

All datasets are chronologically split into training, validation, and
test sets with a ratio of $6:2:2$ and standardized using statistics
computed from the training set. We use 12 historical steps to predict
the subsequent 12 steps. FAST-STG employs UGNet with a hidden dimension of 32, while the
conditional coordinate estimator uses a hidden dimension of 128 and
bounds the predicted scale within $[10^{-3},5]$. The location
correction scale in Eq.~\eqref{eq:conditional_location_scale} is set to
$\alpha_{\mu}=0.1$. The coordinate estimator is pretrained for five
epochs and then frozen before transport learning. FAST-STG is
optimized using Adam with a learning rate of $2\times10^{-3}$ and a
batch size of 8 for at most 300 epochs. Early stopping is applied with
a patience of 10 epochs.

During inference, the future innovation state is initialized from
$\mathcal{N}(\mathbf{0},\mathbf{I})$ and evolved through $K=20$
Euler--Maruyama steps. The numerical stability constant in the
endpoint-induced transport field is set to $\delta=10^{-3}$. We use
$g(\tau)=\gamma(1-\tau)$ with $\gamma=0.08$, $0.08$, and $0.05$ for
AIR-BJ, AIR-GZ, and the PEMS datasets, respectively. Unless otherwise
specified, $S=100$ forecast trajectories are generated for probabilistic
evaluation. All experiments are conducted on a single NVIDIA GeForce
RTX 2080 Ti GPU.

\subsection{Probabilistic Forecasting Performance}
\label{sec:probabilistic_performance}

We evaluate FAST-STG on AIR-BJ, AIR-GZ, and PEMS08.
Table~\ref{tab:probabilistic_results} reports MAE and RMSE for point
forecasting accuracy and CRPS for probabilistic forecast quality.

FAST-STG achieves the lowest CRPS on all three datasets, yielding
relative reductions of approximately $12.1\%$, $5.3\%$, and $16.7\%$
over the best baseline results on AIR-BJ, AIR-GZ, and PEMS08,
respectively. It also achieves the lowest MAE across all three datasets
and the lowest RMSE on AIR-GZ and PEMS08. On AIR-BJ, CW-Diff obtains
the lowest RMSE, while FAST-STG ranks second in RMSE and achieves the
best MAE and CRPS. Overall, these results indicate that FAST-STG
improves probabilistic forecast quality while retaining competitive
point forecasting accuracy.

Figure~\ref{fig:qualitative_results} provides representative
probabilistic forecasts from FAST-STG, DiffSTG, and CoST. In the
illustrated cases, FAST-STG closely tracks the observed future
trajectories while producing relatively concentrated prediction
intervals. The predictive mean remains responsive to local temporal
changes across both air quality and traffic flow scenarios, while the
associated intervals reflect the uncertainty surrounding the future
trajectory. Compared with the baselines, FAST-STG exhibits a favorable
balance between trajectory accuracy and uncertainty concentration in
these representative cases. These examples complement the quantitative
CRPS comparison by illustrating how the learned predictive distribution
behaves across different urban sensing scenarios.

\begin{figure}[h]
    \centering
    \includegraphics[width=\columnwidth]{Figs/Combined.pdf}
    \caption{
    Qualitative comparison on AIR-BJ, AIR-GZ, and PEMS08.
Rows show representative nodes and columns compare FAST-STG, DiffSTG,
and CoST. The vertical dashed line marks the forecast origin; solid
lines, dashed lines, and shaded regions indicate the ground truth,
predictive mean, and 90\% prediction interval.
    }
    \label{fig:qualitative_results}
\end{figure}

\subsection{Component Ablation Study}
\label{sec:ablation}

We conduct controlled ablations on AIR-BJ, AIR-GZ, and PEMS08 to
examine the contributions of the key design choices in FAST-STG.
We consider five variants. \emph{Joint-window raw} performs stochastic
transport over the concatenated historical and future window in the raw
observation space. \emph{Future-only raw} restricts stochastic
evolution to the unknown future, isolating the effect of future-only
transport. \emph{Future-only mean-centered} further removes the
conditional location but does not apply scale normalization.
\emph{Location-scale w/o graph} removes graph conditioning from the
coordinate estimator, while \emph{location-scale w/o freezing} jointly
optimizes the coordinate estimator and transport model instead of
pretraining and freezing the estimator. All variants use the same UGNet
backbone, data splits, optimization budget, and $K=20$ sampling steps.

\begin{figure}[h]
    \centering
    \includegraphics[width=\columnwidth]{Figs/Figure5-消融实验图.pdf}
    \caption{
    Component ablation results on AIR-BJ, AIR-GZ, and PEMS08 in terms
    of MAE, RMSE, and CRPS.
    }
    \label{fig:ablation_visualization}
\end{figure}

Figure~\ref{fig:ablation_visualization} compares the component variants
of FAST-STG. Moving from joint-window transport to future-only transport
generally improves forecasting performance, supporting the benefit of
aligning the stochastic state with the unknown future. Introducing
conditional centering and scale normalization provides further gains,
while removing graph conditioning degrades the learned coordinates.
Jointly optimizing the coordinate estimator and transport model also
leads to consistently worse performance than the staged training
strategy. The complete FAST-STG achieves the lowest MAE, RMSE, and CRPS
across all three datasets, indicating that future-only transport,
graph-conditioned location and scale coordinates, and coordinate
freezing provide complementary benefits.

\subsection{Sensitivity and Efficiency Analysis}
\label{sec:sensitivity}

We analyze the sensitivity of FAST-STG to the number of sampling steps
$K$, stochastic intensity $\gamma$, number of generated trajectories
$S$, and coordinate pretraining epochs $E_{\mathrm{LS}}$ on AIR-GZ.
As shown in Figure~\ref{fig:sensitivity}, FAST-STG achieves the best
overall performance at $K=20$ and $E_{\mathrm{LS}}=5$. Performance
remains stable for $\gamma\in[0.02,0.10]$, with a slight degradation
at $\gamma=0.15$. Increasing $S$ generally improves forecasting
quality: MAE and RMSE improve substantially up to $S=100$ and then
saturate, while CRPS continues to decrease slightly at $S=200$.
Considering both forecasting quality and sampling cost, we adopt
$K=20$, $\gamma=0.08$, $S=100$, and $E_{\mathrm{LS}}=5$ as the
default configuration.

\begin{figure}[h]
    \centering
    \includegraphics[width=\columnwidth]
    {Figs/AIR-GZ_hyperparameters_2x2.pdf}
    \caption{
    Sensitivity analysis of FAST-STG on AIR-GZ with respect to
    (a) the number of sampling steps $K$,
    (b) stochastic intensity $\gamma$,
    (c) the number of generated trajectories $S$, and
    (d) coordinate pretraining epochs $E_{\mathrm{LS}}$.
    MAE and RMSE are shown on the left axis, while CRPS is shown
    on the right axis. Lower values indicate better performance.
    }
    \label{fig:sensitivity}
\end{figure}

\paragraph{Computational Efficiency.}
Table~\ref{tab:efficiency} compares inference efficiency on AIR-BJ
under the same hardware and sampling protocol. FAST-STG contains
$0.910$M parameters and requires $20$ neural function evaluations
(NFEs) per trajectory. Its inference latency is $10.36$ ms per
trajectory, lower than DiffSTG and CoST but higher than the remaining
baselines. FAST-STG also maintains a compact parameter footprint and
moderate GPU memory consumption of $0.24$ GB. Together with the
forecasting results in Table~\ref{tab:probabilistic_results}, these
results characterize the trade-off between probabilistic forecast
quality and iterative sampling cost.

\begin{table}[h]
\centering
\caption{Inference efficiency comparison on AIR-BJ. All methods are
evaluated under the same hardware and sampling protocol.}
\label{tab:efficiency}
\small
\setlength{\tabcolsep}{4pt}
\resizebox{\columnwidth}{!}{
\begin{tabular}{lccccc}
\toprule
Method & Params (M) & NFE & Time (s) & ms/Traj. & Mem. (GB) \\
\midrule
DiffSTG  & 0.840  & 40  & 3416.34 & 19.63 & 2.86 \\
SpecSTG  & 0.073  & 600 & 436.40  & 2.53  & 0.13 \\
NsDiff   & 9.423  & 20  & 682.56  & 3.92  & 0.20 \\
CW-Diff  & 19.403 & 20  & 780.11  & 4.48  & 0.27 \\
CW-Flow  & 10.558 & 20  & 245.23  & 1.41  & 0.30 \\
CoST     & 1.006  & 50  & 1812.33 & 10.48 & 0.02 \\
FAST-STG & 0.910  & 20  & 1802.42 & 10.36 & 0.24 \\
\bottomrule
\end{tabular}}
\end{table}

\section{Conclusion}
\label{sec:conclusion}

This paper proposes FAST-STG, a forecasting-aligned stochastic
transport framework that addresses generation-object and
generation-coordinate misalignment in probabilistic spatio-temporal
graph forecasting. It keeps observed history fixed, restricts stochastic
evolution to the unknown future, and performs transport in graph- and
history-conditioned location-scale coordinates. A staged training strategy further separates coordinate estimation from transport learning. Experiments on six real-world air quality and traffic flow datasets demonstrate improved probabilistic forecast quality with competitive point forecasting accuracy, supported by extensive ablation and analysis.

\bibliographystyle{ACM-Reference-Format}
\bibliography{sample-base}

\clearpage
\appendix

\section{Additional Point Forecasting Results}
\label{app:point_results}

We further evaluate the point forecasting performance of FAST-STG on
PEMS03, PEMS04, PEMS07, and PEMS08. For FAST-STG, point predictions are
obtained by averaging the generated future trajectories. Following the
main experimental protocol, we report MAE, RMSE, and MAPE over the
12-step forecasting horizon. Table~\ref{tab:point_results} compares
FAST-STG with the deterministic spatio-temporal forecasting baselines
introduced in Section~\ref{sec:baselines}, providing an additional
assessment of whether probabilistic generation preserves competitive
point forecasting accuracy.

Table~\ref{tab:point_results} shows that FAST-STG remains competitive
with dedicated deterministic forecasting models across the four
traffic flow datasets. In particular, FAST-STG achieves the second-best
MAE and MAPE on PEMS03 and the second-best RMSE on PEMS08, while
maintaining comparable performance on the remaining datasets. These
results indicate that learning a probabilistic predictive distribution
does not substantially compromise the point forecasting capability of
FAST-STG.

\begin{table*}[t]
\centering
\caption{
Point forecasting performance on PEMS03, PEMS04, PEMS07, and PEMS08. Best results are shown in
bold and second-best results are underlined.
}
\label{tab:point_results}
\small
\setlength{\tabcolsep}{3.6pt}
\renewcommand{\arraystretch}{1.08}
\resizebox{\textwidth}{!}{
\begin{tabular}{@{}lcccccccccccc@{}}
\toprule
&
\multicolumn{3}{c}{PEMS03} &
\multicolumn{3}{c}{PEMS04} &
\multicolumn{3}{c}{PEMS07} &
\multicolumn{3}{c}{PEMS08} \\
\cmidrule(lr){2-4}
\cmidrule(lr){5-7}
\cmidrule(lr){8-10}
\cmidrule(lr){11-13}

Method
& MAE $\downarrow$ & RMSE $\downarrow$ & MAPE (\%) $\downarrow$
& MAE $\downarrow$ & RMSE $\downarrow$ & MAPE (\%) $\downarrow$
& MAE $\downarrow$ & RMSE $\downarrow$ & MAPE (\%) $\downarrow$
& MAE $\downarrow$ & RMSE $\downarrow$ & MAPE (\%) $\downarrow$ \\
\midrule

STGCN~\cite{yu2018spatio}
& 17.49 & 30.12 & 17.15
& 22.70 & 35.55 & 14.59
& 25.38 & 38.78 & 11.08
& 18.02 & 27.83 & 11.40 \\

DCRNN~\cite{li2018diffusion}
& 18.18 & 30.31 & 18.91
& 24.70 & 38.12 & 17.12
& 25.30 & 38.58 & 11.66
& 17.86 & 27.83 & 11.45 \\

Graph WaveNet~\cite{wu2019graph}
& 19.85 & 32.94 & 19.31
& 25.45 & 39.70 & 17.29
& 26.85 & 42.78 & 12.12
& 19.13 & 31.05 & 12.68 \\

ASTGCN~\cite{guo2019attention}
& 17.69 & 29.66 & 19.40
& 22.93 & 35.22 & 16.56
& 28.05 & 42.57 & 13.92
& 18.61 & 28.16 & 13.08 \\

AGCRN~\cite{bai2020adaptive}
& \textbf{15.24} & \textbf{26.65} & 15.89
& \underline{19.38} & \underline{31.25} & \underline{13.40}
& \underline{20.57} & \underline{34.40} & \textbf{8.74}
& \textbf{15.32} & \textbf{24.41} & \underline{10.03} \\

STGODE~\cite{fang2021stgode}
& 16.50 & 27.84 & 16.69
& 20.84 & 32.82 & 13.77
& 22.99 & 37.54 & 10.14
& 16.81 & 25.97 & 10.62 \\

STG-NCDE~\cite{choi2022stgncde}
& 15.57 & \underline{27.09} & \textbf{15.06}
& \textbf{19.21} & \textbf{31.09} & \textbf{12.76}
& \textbf{20.53} & \textbf{33.84} & \underline{8.80}
& \underline{15.45} & 24.81 & \textbf{9.92} \\

\midrule

\textbf{FAST-STG}
& \underline{15.49} & 29.87 & \underline{15.84}
& 20.11 & 32.08 & 13.49
& 22.15 & 34.67 & 9.17
& 15.66 & \underline{24.53} & 10.25 \\

\bottomrule
\end{tabular}}
\end{table*}

\section{Additional Experimental Details}
\label{app:experimental_details}

This section provides supplementary details to improve the
reproducibility of FAST-STG. We first summarize the two-stage training
procedure and stochastic inference process in algorithmic form, and
then report the exact numerical results underlying the sensitivity
analysis presented in the main text. These additional details
complement the methodological descriptions in Section~\ref{sec:method}
and the experimental settings in Section~\ref{sec:experimental_setup}.

\subsection{Training Procedure}
\label{app:training_procedure}

Algorithm~\ref{alg:training} summarizes the complete two-stage training
procedure of FAST-STG. In Stage I, the conditional coordinate estimator
$f_{\phi}$ is optimized using the location-scale objective in
Eq.~\eqref{eq:ls_loss}. The best estimator parameters
$\phi^{\star}$ are selected according to the validation negative
log-likelihood and then frozen. In Stage II, the fixed coordinate
estimator defines standardized future innovations, and the transport
backbone $\operatorname{UGNet}_{\theta}$ is optimized using the endpoint
prediction objective in Eq.~\eqref{eq:transport_loss}. Model selection
is performed according to validation CRPS with early stopping. This
staged procedure keeps the conditional coordinate system fixed
throughout transport learning.

\begin{algorithm}[h]
\caption{Two-Stage Training of FAST-STG}
\label{alg:training}
\begin{algorithmic}[1]

\Require Training set $\mathcal{D}_{\mathrm{tr}}$;
validation set $\mathcal{D}_{\mathrm{val}}$;
spatial graph $\mathbf{A}$;
coordinate estimator $f_{\phi}$;
transport backbone $\operatorname{UGNet}_{\theta}$;
maximum epoch budget $E$;
coordinate pretraining epochs $E_{\mathrm{LS}}$

\Ensure Selected parameters $\phi^{\star}$ and $\theta^{\star}$

\Statex
\Statex \textbf{Stage I: Conditional coordinate pretraining}

\State Freeze $\theta$ and enable gradients for $\phi$
\State $\ell_{\mathrm{LS}}^{\star}\gets+\infty$

\For{$e=1,\ldots,E_{\mathrm{LS}}$}

    \ForAll{minibatches
    $(\mathbf{X}^{\mathrm{his}},
    \mathbf{X}^{\mathrm{fut}})
    \sim\mathcal{D}_{\mathrm{tr}}$}

        \State $(\boldsymbol{\mu}_{\phi},
        \boldsymbol{\sigma}_{\phi})
        \gets
        f_{\phi}
        (\mathbf{X}^{\mathrm{his}},\mathbf{A})$

        \State Compute $\mathcal{L}_{\mathrm{LS}}(\phi)$
        using Eq.~\eqref{eq:ls_loss}

        \State Update $\phi$ by minimizing
        $\mathcal{L}_{\mathrm{LS}}$

    \EndFor

    \State $\ell_{\mathrm{LS}}
    \gets
    \operatorname{ValNLL}
    (f_{\phi},\mathcal{D}_{\mathrm{val}})$

    \If{$\ell_{\mathrm{LS}}
    <\ell_{\mathrm{LS}}^{\star}$}

        \State $\phi^{\star}\gets\phi$

        \State $\ell_{\mathrm{LS}}^{\star}
        \gets\ell_{\mathrm{LS}}$

    \EndIf

\EndFor

\State Restore $\phi\gets\phi^{\star}$
\State Freeze $\phi$ and enable gradients for $\theta$

\Statex
\Statex \textbf{Stage II: Future-only transport learning}

\For{$e=E_{\mathrm{LS}}+1,\ldots,E$}

    \ForAll{minibatches
    $(\mathbf{X}^{\mathrm{his}},
    \mathbf{X}^{\mathrm{fut}})
    \sim\mathcal{D}_{\mathrm{tr}}$}

        \State $(\boldsymbol{\mu}_{\phi^{\star}},
        \boldsymbol{\sigma}_{\phi^{\star}})
        \gets
        f_{\phi^{\star}}
        (\mathbf{X}^{\mathrm{his}},\mathbf{A})$

        \State $\mathbf{R}^{\mathrm{fut}}
        \gets
        \operatorname{sg}
        \left[
        \left(
        \mathbf{X}^{\mathrm{fut}}
        -
        \boldsymbol{\mu}_{\phi^{\star}}
        \right)
        \oslash
        \boldsymbol{\sigma}_{\phi^{\star}}
        \right]$

        \State Sample
        $\mathbf{Z}_{0}
        \sim
        \mathcal{N}(\mathbf{0},\mathbf{I})$
        and
        $\tau\sim\mathcal{U}(0,1)$

        \State $\mathbf{Z}_{\tau}
        \gets
        (1-\tau)\mathbf{Z}_{0}
        +
        \tau\mathbf{R}^{\mathrm{fut}}$

        \State $\overline{\mathbf{Z}}_{\tau}
        \gets
        \operatorname{Cat}_{t}
        \left(
        \mathbf{X}^{\mathrm{his}},
        \mathbf{Z}_{\tau}
        \right)$

        \State $\mathbf{C}
        \gets
        \operatorname{Cat}_{t}
        \left(
        \mathbf{X}^{\mathrm{his}},
        \mathbf{0}
        \right)$

        \State $\widehat{\mathbf{R}}_{\theta}^{\mathrm{fut}}
        \gets
        \Pi_{\mathrm{fut}}
        \left[
        \operatorname{UGNet}_{\theta}
        \left(
        \overline{\mathbf{Z}}_{\tau},
        \tau;
        \mathbf{C},
        \mathbf{A}
        \right)
        \right]$

        \State Compute
        $\mathcal{L}_{\mathrm{T}}(\theta)$
        using Eq.~\eqref{eq:transport_loss}

        \State Update $\theta$ while keeping
        $\phi^{\star}$ frozen

    \EndFor

    \State Evaluate validation CRPS using a fixed
    number of forecast trajectories

    \If{validation CRPS improves}
        \State $\theta^{\star}\gets\theta$
    \EndIf

    \If{early-stopping criterion is satisfied}
        \State \textbf{break}
    \EndIf

\EndFor

\State \Return $\phi^{\star},\theta^{\star}$

\end{algorithmic}
\end{algorithm}

\subsection{Inference Procedure}
\label{app:inference_procedure}

Algorithm~\ref{alg:inference} summarizes the complete stochastic
inference procedure of FAST-STG. The selected coordinate estimator
$f_{\phi^{\star}}$ and transport backbone
$\operatorname{UGNet}_{\theta^{\star}}$ remain fixed throughout
inference. The conditional location and scale are first computed from
the normalized historical observations and spatial graph. For each
forecast trajectory, FAST-STG initializes the future innovation state
from a standard Gaussian distribution and iteratively predicts the
clean future endpoint from the current state. The corresponding
endpoint-induced transport direction is combined with a time-dependent
stochastic term and integrated using Euler--Maruyama updates. After
$K$ steps, the generated innovation is transformed back through the
conditional location and scale coordinates. Repeating the procedure
produces $S$ probabilistic forecast trajectories.

\begin{algorithm}[h]
\caption{Stochastic Inference of FAST-STG}
\label{alg:inference}
\begin{algorithmic}[1]

\Require Normalized observed history
$\mathbf{X}^{\mathrm{his}}$;
spatial graph $\mathbf{A}$;
selected coordinate estimator $f_{\phi^{\star}}$;
selected transport backbone
$\operatorname{UGNet}_{\theta^{\star}}$;
sampling steps $K$;
stochastic intensity $\gamma$;
number of trajectories $S$;
stability constant $\delta$

\Ensure Forecast trajectories
$\left\{
\widehat{\mathbf{X}}^{\mathrm{fut},(s)}
\right\}_{s=1}^{S}$

\State $(\boldsymbol{\mu}_{\phi^{\star}},
\boldsymbol{\sigma}_{\phi^{\star}})
\gets
f_{\phi^{\star}}
\left(
\mathbf{X}^{\mathrm{his}},
\mathbf{A}
\right)$

\State $\mathbf{C}
\gets
\operatorname{Cat}_{t}
\left(
\mathbf{X}^{\mathrm{his}},
\mathbf{0}
\right)$

\State $\Delta\tau\gets 1/K$

\For{$s=1,\ldots,S$}

    \State Sample
    $\mathbf{Z}_{0}^{(s)}
    \sim
    \mathcal{N}(\mathbf{0},\mathbf{I})$

    \For{$k=0,\ldots,K-1$}

        \State $\tau_k\gets k\Delta\tau$

        \State $\overline{\mathbf{Z}}_{k}^{(s)}
        \gets
        \operatorname{Cat}_{t}
        \left(
        \mathbf{X}^{\mathrm{his}},
        \mathbf{Z}_{k}^{(s)}
        \right)$

        \State $\widehat{\mathbf{R}}_{\theta^{\star},k}^{\mathrm{fut},(s)}
        \gets
        \Pi_{\mathrm{fut}}
        \left[
        \operatorname{UGNet}_{\theta^{\star}}
        \left(
        \overline{\mathbf{Z}}_{k}^{(s)},
        \tau_k;
        \mathbf{C},
        \mathbf{A}
        \right)
        \right]$

        \State $\mathbf{v}_{k}^{(s)}
        \gets
        \dfrac{
        \widehat{\mathbf{R}}_{\theta^{\star},k}^{\mathrm{fut},(s)}
        -
        \mathbf{Z}_{k}^{(s)}
        }{
        1-\tau_k+\delta
        }$

        \State $g_k\gets\gamma(1-\tau_k)$

        \State Sample
        $\boldsymbol{\epsilon}_{k}^{(s)}
        \sim
        \mathcal{N}(\mathbf{0},\mathbf{I})$

        \State $\mathbf{Z}_{k+1}^{(s)}
        \gets
        \mathbf{Z}_{k}^{(s)}
        +
        \Delta\tau\,
        \mathbf{v}_{k}^{(s)}
        +
        g_k\sqrt{\Delta\tau}\,
        \boldsymbol{\epsilon}_{k}^{(s)}$

    \EndFor

    \State $\widehat{\mathbf{X}}^{\mathrm{fut},(s)}
    \gets
    \boldsymbol{\mu}_{\phi^{\star}}
    +
    \boldsymbol{\sigma}_{\phi^{\star}}
    \odot
    \mathbf{Z}_{K}^{(s)}$

\EndFor

\State \Return
$\left\{
\widehat{\mathbf{X}}^{\mathrm{fut},(s)}
\right\}_{s=1}^{S}$

\end{algorithmic}
\end{algorithm}

\subsection{Detailed Sensitivity Results}
\label{app:sensitivity_details}

To complement the compact sensitivity visualization in
Figure~\ref{fig:sensitivity}, we provide enlarged results for each
hyperparameter on AIR-GZ. In each experiment, only the parameter under
study is varied while the remaining settings are fixed to their default
values. Figures~\ref{fig:sensitivity_k}--\ref{fig:sensitivity_els}
report the corresponding MAE, RMSE, and CRPS values.

\begin{figure}[t]
    \centering
    \includegraphics[width=\columnwidth]
    {Figs/AIR-GZ_K_sensitivity.pdf}
    \caption{
    Sensitivity to the number of sampling steps $K$ on AIR-GZ.
    FAST-STG achieves the best overall performance at $K=20$.
    Lower values are better for all metrics.
    }
    \label{fig:sensitivity_k}
\end{figure}

\begin{figure}[t]
    \centering
    \includegraphics[width=\columnwidth]
    {Figs/AIR-GZ_gamma_sensitivity.pdf}
    \caption{
    Sensitivity to the stochastic intensity $\gamma$ on AIR-GZ.
    Performance remains stable for $\gamma\in[0.02,0.10]$ and
    slightly degrades at $\gamma=0.15$.
    Lower values are better for all metrics.
    }
    \label{fig:sensitivity_gamma}
\end{figure}

\begin{figure}[t]
    \centering
    \includegraphics[width=\columnwidth]
    {Figs/AIR-GZ_S_sensitivity.pdf}
    \caption{
    Sensitivity to the number of generated trajectories $S$ on AIR-GZ.
    MAE and RMSE attain their lowest values at $S=100$, while CRPS
    continues to improve at $S=200$.
    Lower values are better for all metrics.
    }
    \label{fig:sensitivity_s}
\end{figure}

\begin{figure}[t]
    \centering
    \includegraphics[width=\columnwidth]
    {Figs/AIR-GZ_ELS_sensitivity.pdf}
    \caption{
    Sensitivity to the coordinate pretraining epochs
    $E_{\mathrm{LS}}$ on AIR-GZ.
    The best overall performance is obtained with
    $E_{\mathrm{LS}}=5$.
    Lower values are better for all metrics.
    }
    \label{fig:sensitivity_els}
\end{figure}

For the number of sampling steps, $K=20$ achieves an MAE of
$10.0037$, an RMSE of $15.0019$, and a CRPS of $0.1821$, outperforming
the other evaluated step counts. The results are comparatively
insensitive to the stochastic intensity: $\gamma=0.02$, $0.05$,
$0.08$, and $0.10$ produce nearly identical performance, whereas
$\gamma=0.15$ results in a slight degradation. We therefore use
$\gamma=0.08$ as the default setting.

Increasing the number of generated trajectories generally improves
forecast quality. At $S=100$, FAST-STG obtains its lowest MAE and RMSE
of $10.0037$ and $15.0019$, respectively, while increasing to $S=200$
further reduces CRPS from $0.1821$ to $0.1773$ with little change in
point accuracy. Considering forecasting quality and sampling cost, we
use $S=100$ by default. Finally, coordinate pretraining performs best
at $E_{\mathrm{LS}}=5$, yielding an MAE of $10.0037$, an RMSE of
$15.0019$, and a CRPS of $0.1821$. These results support the default
configuration of $K=20$, $\gamma=0.08$, $S=100$, and
$E_{\mathrm{LS}}=5$ used in the main experiments.

\section{Additional Probabilistic Evaluation}
\label{app:interval_metrics}

In addition to CRPS reported in the main text, we further evaluate the
quality of predictive intervals using three complementary metrics:
prediction interval coverage probability (PICP), prediction interval
width (PIW), and mean interval score (MIS). These metrics characterize
different aspects of the predictive distribution, including empirical
coverage, interval sharpness, and the joint quality of interval width
and coverage. We report all three metrics at the 90\% prediction level on AIR-BJ,
providing a complementary assessment of interval calibration and
sharpness under the unified evaluation setting.

\subsection{Interval-Based Probabilistic Metrics}
\label{app:interval_metric_definitions}

Given $S$ generated forecast trajectories, we construct a 90\%
prediction interval for each target entry from the empirical 5th and
95th percentiles. Let $y_i$ denote the ground-truth value and
$l_i$ and $u_i$ denote the corresponding lower and upper bounds,
where $i=1,\ldots,M$ indexes all evaluated future entries.

The prediction interval coverage probability (PICP) measures the
empirical coverage of the prediction intervals:
\begin{equation}
\operatorname{PICP}_{90}
=
\frac{1}{M}
\sum_{i=1}^{M}
\mathbb{I}
\left(
l_i \leq y_i \leq u_i
\right),
\label{eq:picp90}
\end{equation}
where $\mathbb{I}(\cdot)$ is the indicator function. For a nominal
90\% interval, values close to $0.90$ indicate agreement between
empirical and nominal coverage.

The prediction interval width (PIW) evaluates interval sharpness:
\begin{equation}
\operatorname{PIW}_{90}
=
\frac{1}{M}
\sum_{i=1}^{M}
(u_i-l_i).
\label{eq:piw90}
\end{equation}
Under comparable coverage, a smaller PIW indicates sharper predictive
intervals.

We further use the mean interval score (MIS) to jointly evaluate
interval width and coverage errors. For a $(1-\alpha)$ prediction
interval, the interval score is
\begin{equation}
\begin{aligned}
\operatorname{IS}_{\alpha,i}
={}&
(u_i-l_i)
+
\frac{2}{\alpha}(l_i-y_i)
\mathbb{I}(y_i<l_i)
\\
&+
\frac{2}{\alpha}(y_i-u_i)
\mathbb{I}(y_i>u_i),
\end{aligned}
\label{eq:interval_score}
\end{equation}
and for the 90\% interval, with $\alpha=0.1$,
\begin{equation}
\operatorname{MIS}_{90}
=
\frac{1}{M}
\sum_{i=1}^{M}
\operatorname{IS}_{0.1,i}.
\label{eq:mis90}
\end{equation}
A lower MIS indicates better overall interval quality by jointly
penalizing excessive width and observations falling outside the
predicted interval.

\subsection{Interval Forecasting Results}
\label{app:interval_results}

Table~\ref{tab:interval_results} reports the 90\% interval forecasting
results on AIR-BJ. All methods are evaluated using the same data split,
preprocessing procedure, and number of generated trajectories. PICP
measures empirical coverage with a nominal target of $0.90$, while PIW
reflects interval sharpness under comparable coverage. MIS jointly
accounts for interval width and coverage errors, with lower values
indicating better overall interval quality.

\begin{table}[h]
\centering
\caption{
Interval-based probabilistic forecasting results on AIR-BJ.
PICP is evaluated against the nominal coverage level of $0.90$.
Lower PIW indicates sharper intervals under comparable coverage,
while lower MIS indicates better overall interval quality.
}
\label{tab:interval_results}
\setlength{\tabcolsep}{6pt}
\renewcommand{\arraystretch}{1.08}
\begin{tabular}{lccc}
\toprule
Method & PICP$_{90}$ & PIW$_{90}$ & MIS$_{90}$ \\
\midrule
DiffSTG  & 0.6987 & 56.54 & 157.91 \\
SpecSTG  & 0.6415 & 44.00 & 182.73 \\
NsDiff   & 0.8403 & 53.82 & 144.04 \\
CW-Diff  & 0.4221 & 24.20 & 281.91 \\
CW-Flow  & 0.3397 & 13.28 & 265.15 \\
CoST     & \underline{0.8693} & 60.83 & \underline{108.12} \\
\midrule
FAST-STG & \textbf{0.9014} & 69.18 & \textbf{102.38} \\
\bottomrule
\end{tabular}
\end{table}

FAST-STG achieves a PICP of $0.9014$, closely matching the nominal
coverage level of $0.90$, while CoST provides the next closest
coverage at $0.8693$. Although FAST-STG produces wider prediction
intervals, the substantially lower coverage of several competing
methods indicates that their narrower intervals are associated with
under-coverage. FAST-STG also obtains the lowest MIS of $102.38$,
improving upon CoST by approximately $5.3\%$. These results indicate
that FAST-STG provides a favorable balance between empirical coverage
and interval sharpness, complementing the CRPS results reported in the
main text.

\end{document}
\endinput
%%
%% End of file `sigconf.tex'.
